\documentclass[11pt,a4paper]{article}
\usepackage[T1]{fontenc}
\usepackage[utf8]{inputenc}
\usepackage[provide=*,lithuanian]{babel}
\usepackage{lmodern,amsmath,amssymb,enumitem}
\usepackage[margin=1.5cm]{geometry}
\setlength{\parindent}{0pt}
\setlist[enumerate]{label=\textbf{\arabic*.},leftmargin=0.65cm,itemsep=0.28cm,topsep=0.2cm,parsep=0pt}
\setlength{\abovedisplayskip}{4pt}
\setlength{\belowdisplayskip}{3pt}
\setlength{\abovedisplayshortskip}{2pt}
\setlength{\belowdisplayshortskip}{3pt}
\begin{document}
\begin{center}
{\Large\bfseries Sumos ir skirtumo kvadrato\\[2pt]formulių taikymas}

\medskip
{\large 50 mišrių uždavinių}
\end{center}
\textbf{Užduotis.} Išspręskite uždavinius, taikydami sumos ir skirtumo
kvadrato formules.
\begin{enumerate}
\item Apskaičiuokite kuo trumpesniu būdu:\[1203^2-2\cdot1203\cdot1198+1198^2.\]
\item Skaidykite dauginamaisiais, atpažindami pilnąjį kvadratą:\[16+24ab+9a^2b^2.\]
\item Neapskaičiuodami \(997^2\), raskite reikšmę:\[997^2+6\cdot997+9.\]
\item Žinoma, kad \(a+b=11\) ir \(ab=24\). Nerasdami atskirai \(a\) ir \(b\), raskite \(a^2+b^2\).
\item Neatlikdami triženklių skaičių daugybos, apskaičiuokite:\[\sqrt{601^2+2\cdot601\cdot399+399^2}.\]
\item Suprastinkite reiškinį ir nurodykite, su kuriomis kintamųjų reikšmėmis pradinis reiškinys turi prasmę:\[\frac{4x^2+12xy+9y^2}{2x+3y}.\]
\item Skaidykite dauginamaisiais, atpažindami pilnąjį kvadratą:\[\frac{x^2}{4}+xy+y^2.\]
\item Žinoma, kad \(x^2+y^2=50\), \(xy=7\) ir \(x>y\). Raskite \(x-y\).
\item Skaidykite dauginamaisiais, atpažindami pilnąjį kvadratą:\[3x^2-12xy+12y^2.\]
\item Išspręskite lygtį nenaudodami diskriminanto:\[4x^2+12x+9=25.\]
\item Įrodykite tapatybę, naudodami sumos ir skirtumo kvadrato formules:\[(a+b)^2=(a-b)^2+4ab.\]
\item Žinoma, kad \(a^2+b^2=34\), \(ab=15\) ir \(a+b>0\). Raskite \(a+b\).
\item Įrodykite nelygybę visiems realiesiems \(a,b\) ir nustatykite lygybės atvejį:\[a^2+b^2\ge2ab.\]
\item Įrodykite tapatybę, naudodami sumos ir skirtumo kvadrato formules:\[(a+2b)^2+(a-2b)^2=2a^2+8b^2.\]
\item Atpažinę pilnąjį kvadratą, suprastinkite:\[\frac{1001^2-2\cdot1001+1}{500^2}.\]
\item Žinoma, kad \(a-b=7\) ir \(ab=-6\). Raskite \(a^2+b^2\).
\item Išspręskite lygtį nenaudodami diskriminanto:\[x^2-10x+25=3.\]
\item Skaidykite dauginamaisiais, atpažindami pilnąjį kvadratą:\[9a^4-30a^2b+25b^2.\]
\item Dviejų skaičių suma yra \(5\), o jų kvadratų suma -- \(13\). Raskite šių skaičių sandaugą.
\item Raskite mažiausią galimą reiškinio reikšmę ir nurodykite, kada ji gaunama (kintamieji realieji):\[3x^2-12x+17.\]
\item Įrodykite nelygybę visiems \(x\in\mathbb R\) ir nustatykite lygybės atvejus:\[4x^2-12x+11\ge2.\]
\item Raskite visas realiųjų skaičių poras \((x,y)\), kurioms:\[x^2+4y^2-6x+8y+13=0.\]
\item Išspręskite lygtį nenaudodami diskriminanto:\[9x^2-6x+1=2.\]
\item Raskite mažiausią galimą reiškinio reikšmę ir nurodykite, kada ji gaunama (kintamieji realieji):\[4x^2+4x+7.\]
\item Suprastinkite reiškinį ir nurodykite, su kuriomis kintamųjų reikšmėmis pradinis reiškinys turi prasmę:\[\frac{x^2-6x+9}{3-x}.\]
\item Skaidykite dauginamaisiais, atpažindami pilnąjį kvadratą:\[(x-1)^2+6(x-1)+9.\]
\item Raskite mažiausią galimą reiškinio reikšmę ir nurodykite, kada ji gaunama (kintamieji realieji):\[x^2+6x+12.\]
\item Neapskaičiuodami kvadratų įprastu būdu, apskaičiuokite:\[1003^2+997^2-2\cdot1000^2.\]
\item Išspręskite lygtį pasinaudoję greitąja daugyba:\[(x-2)^2+4(x-2)+4=49.\]
\item Suprastinkite reiškinį ir nurodykite, su kuriomis kintamųjų reikšmėmis pradinis reiškinys turi prasmę:\[\frac{x^2+2x+1}{x^2-1}.\]
\item Išspręskite lygtį nenaudodami diskriminanto:\[x^2+6x+5=0.\]
\item Kvadrato kraštinę pailginus \(3\) cm, jo plotas padidėjo \(69\text{ cm}^2\). Raskite pradinio kvadrato kraštinės ilgį.
\item Įrodykite nelygybę visiems \(x\in\mathbb R\) ir nustatykite, kada galioja lygybė:\[x^2+6x+10\ge1.\]
\item Stačiakampio perimetras yra \(28\) cm, o plotas -- \(40\text{ cm}^2\). Raskite įstrižainės ilgį, nerasdami atskirai kraštinių ilgių.
\item Įrodykite tapatybę, naudodami sumos ir skirtumo kvadrato formules (\(x\ne0\)):\[(x+\frac1x)^2+(x-\frac1x)^2=2x^2+\frac2{x^2}.\]
\item Įrodykite nelygybę visiems realiesiems \(x,y\) ir nustatykite lygybės atvejį:\[x^2+y^2+2x-4y+5\ge0.\]
\item Raskite visas realiųjų skaičių poras \((x,y)\), kurioms:\[9x^2+4y^2-12x+4y+5=0.\]
\item Stačiojo trikampio statinių ilgių suma yra \(17\) cm, o jo plotas -- \(30\text{ cm}^2\). Raskite įžambinės ilgį, nerasdami atskirai statinių ilgių.
\item Suprastinkite reiškinį ir nurodykite, su kuriomis kintamųjų reikšmėmis pradinis reiškinys turi prasmę:\[\frac{a^2-4ab+4b^2}{a^2-2ab}.\]
\item Įrodykite tapatybę, naudodami sumos ir skirtumo kvadrato formules:\[(a+b)^2+(b+c)^2+(c+a)^2=2(a^2+b^2+c^2+ab+bc+ca).\]
\item Raskite visas realiųjų skaičių poras \((x,y)\), kurioms:\[4x^2+y^2+4x-6y+10=0.\]
\item Raskite mažiausią galimą reiškinio reikšmę ir nurodykite, kada ji gaunama (kintamieji realieji):\[a^2+b^2-4a+6b+20.\]
\item Stačiakampio kraštinių ilgių skirtumas yra \(4\) cm, o plotas -- \(45\text{ cm}^2\). Raskite įstrižainės ilgį, nerasdami atskirai kraštinių ilgių.
\item Suprastinkite reiškinį ir nurodykite, su kuriomis kintamųjų reikšmėmis pradinis reiškinys turi prasmę:\[\frac{x^2+4x+4}{x^2-4}:\frac{x+2}{x-2}.\]
\item Raskite visas realiųjų skaičių poras \((x,y)\), kurioms:\[2x^2+3y^2-8x+6y+11=0.\]
\item Įrodykite, kad visiems realiesiems \(x\):\[x^2+x+1>0.\]
\item Raskite didžiausią galimą reiškinio reikšmę ir nurodykite, kada ji gaunama (kintamieji realieji):\[-2x^2+8x+1.\]
\item Dviejų kvadratų plotų suma lygi \(100\text{ cm}^2\), o jų kraštinių ilgių suma -- \(14\) cm. Raskite didesniojo ir mažesniojo kvadratų kraštinių ilgių skirtumą.
\item Įrodykite tapatybę, naudodami sumos ir skirtumo kvadrato formules:\[(a+b)^2+(c+d)^2-(a+d)^2-(b+c)^2=2(a-c)(b-d).\]
\item Raskite visas realiųjų skaičių poras \((x,y)\), kurioms:\[(x+y)^2+(y-3)^2=4(x+y)-4.\]
\end{enumerate}
\clearpage
\section*{Atsakymai}
\begin{enumerate}
\item \(25\).
\item \((3ab+4)^2\).
\item \(1\,000\,000\).
\item \(73\).
\item \(1000\).
\item \(2x+3y\), kai \(2x+3y\ne0\).
\item \((\frac x2+y)^2\).
\item \(6\).
\item \(3(x-2y)^2\).
\item \(x=-4\) arba \(x=1\).
\item Tapatybė teisinga.
\item \(8\).
\item Lygybė, kai \(a=b\).
\item Tapatybė teisinga.
\item \(4\).
\item \(37\).
\item \(x=5\pm\sqrt3\).
\item \((3a^2-5b)^2\).
\item \(6\).
\item \(5\), kai \(x=2\).
\item Lygybė, kai \(x=\frac32\).
\item \((x,y)=(3,-1)\).
\item \(x=\dfrac{1\pm\sqrt2}{3}\).
\item \(6\), kai \(x=-\frac12\).
\item \(3-x\), kai \(x\ne3\).
\item \((x+2)^2\).
\item \(3\), kai \(x=-3\).
\item \(18\).
\item \(x=-7\) arba \(x=7\).
\item \(\dfrac{x+1}{x-1}\), kai \(x\ne\pm1\).
\item \(x=-5\) arba \(x=-1\).
\item \(10\text{ cm}\).
\item Lygybė, kai \(x=-3\).
\item \(2\sqrt{29}\text{ cm}\).
\item Tapatybė teisinga.
\item Lygybė, kai \(x=-1\), \(y=2\).
\item \((x,y)=(\frac23,-\frac12)\).
\item \(13\text{ cm}\).
\item \(\dfrac{a-2b}{a}\), kai \(a\ne0\) ir \(a\ne2b\).
\item Tapatybė teisinga.
\item \((x,y)=(-\frac12,3)\).
\item \(7\), kai \(a=2\), \(b=-3\).
\item \(\sqrt{106}\text{ cm}\).
\item \(1\), kai \(x\ne\pm2\).
\item \((x,y)=(2,-1)\).
\item Nelygybė teisinga visiems realiesiems \(x\).
\item \(9\), kai \(x=2\).
\item \(2\text{ cm}\).
\item Tapatybė teisinga.
\item \((x,y)=(-1,3)\).
\end{enumerate}
\end{document}
